\documentclass[12pt]{article}
\usepackage{amsfonts,amsthm,amsmath,amssymb,mathtools}
\usepackage{mathrsfs}
\usepackage{enumitem}
\usepackage{tikz-cd}
\usepackage{geometry}
\usepackage{mlmodern}

\geometry{letterpaper, portrait, margin=1in}

\setcounter{section}{0}

\renewcommand{\thesection}{\S\arabic{section}}
\renewcommand{\thesubsection}{\arabic{section}}
\newcommand{\dd}{\mathop{}\!\mathrm{d}}
\newcommand{\R}{\mathbb{R}}
\newcommand{\Z}{\mathbb{Z}}
\newcommand{\C}{\mathbb{C}}
\newcommand{\Q}{\mathbb{Q}}

\DeclareMathOperator{\lcm}{lcm}
\DeclareMathOperator{\im}{im}

\newtheorem{thm}{Theorem}
\newenvironment{theorem}[1]{
	\renewcommand{\thethm}{#1}
	\begin{thm}
		}{\end{thm}}

\newtheorem{lma}{Lemma}
\newenvironment{lemma}[1]{
	\renewcommand{\thelma}{#1}
	\begin{lma}
		}{\end{lma}}

\theoremstyle{definition}
\newtheorem{ex}{}
\newenvironment{exercise}[1]{
	\renewcommand{\theex}{#1}
	\begin{ex}
		}{\end{ex}}

\title{Project 2: Jacobians}
\author{Connor Mooneyhan}
\date{December 9, 2025}

\begin{document}

\maketitle

Prompt: The key property that allows us to define a group structure on a cubic curve is Bezout’s theorem that indicates that there are precisely three intersection points of a cubic curve with a line.
The goal of this project is to learn a little bit about higher-dimensional abelian varieties which give a way to associated a group structure with points on other sorts of curves.
Research and write a report that addresses the following qustions: (a) What are divisors and principal divisors on a curve, (b) What is the Jacobian of a curve and how is its group structure defined? (c) How do you explicitly compute with the Jacobian of a genus 2 curve? (d) How is an elliptic curve a Jacobian? and (e) What can be said about the Mordell-Weil group of a Jacobian?

\section{Introduction}

Note: I will attempt to remain rigorous while supressing mention of ``varieties'' and the fields over which they may be defined in favor of the language of ``curves'' in an attempt to demonstrate understanding of the material without hiding behind concepts I don't fully understand the significance of.
\section{Divisors and principal divisors}
Let $C$ be a projective curve.
Then a \textit{divisor} of $C$ is a finite formal sum \begin{equation}\label{divisor}
	\sum_{i=1}^m n_iP_i,
\end{equation}
where each $P_i$ is a point on $C$ and each $n_i \in \mathbb{Z}$.
In other words, the set of divisors is the free abelian group generated by the points of $C$.
This group is called the \textit{divisor group} of $C$ (denoted $\mathrm{Div}(C)$), and addition is defined as usual for linear combinations, where the coefficients of ``like terms'' (i.e. terms corresponding to the same point) are added in the final sum.

We now need the concept of the order of a function, whose definition seems to get further into the weeds of valuation rings than the rest of this project will require.
As such, I will include a couple of theorems about the order $\mathrm{ord}_P(f)$ of a function $f$ at a point $P$ which establish the properties we will need:
\begin{enumerate}
	\item If $\mathrm{ord}_P(f) > 0$, then $f$ has a zero at $P$.
	      If $\mathrm{ord}_P(f) < 0$, then $f$ has a pole (or, is infinite) at $P$.
	      If $\mathrm{ord}_P(f) = 0$, then $f$ is defined but nonzero at $P$.
	\item Given a smooth curve $C$ and a rational function $f$ defined over $C$, there are only finitely many points of $C$ at which $f$ has a pole or a zero.
	      Thus $\mathrm{ord}_P(f) = 0$ at all but finitely many points $P \in C$.
\end{enumerate}
Given that $C$ is smooth, this allows us to assign to any rational function $f$ over $C$ a divisor $\mathrm{div}(f)$ defined by \begin{equation*}
	\mathrm{div}(f) = \sum_{P \in C} \mathrm{ord}_P(f)\cdot P.
\end{equation*}
We can now define a \textit{principal divisor} to be any divisor $D \in \mathrm{Div}(C)$ such that $D = \mathrm{div}(f)$ for some rational function $f$ defined over $C$.

\section{The Jacobian of a curve}
Two divisors $D_1, D_2 \in \mathrm{Div}(C)$ are said to be \textit{linearly equivalent}, denoted $D_1 \sim D_2$, if $D_1 - D_2$ is principal.
We define the Picard group $\mathrm{Pic}(C)$ to be the quotient of $\mathrm{Div}(C)$ by the subgroup of principal divisors, i.e. divisors modulo linear equivalence.
The \textit{degree} of a divisor is the sum of its coefficients: \begin{equation}
	\mathrm{deg}\left( \sum_{i=1}^m n_iP_i \right) = \sum_{i=1}^m n_i.
\end{equation}
We can then define the groups $\mathrm{Div}^0(C)$ and $\mathrm{Pic}^0(C)$ to be the subgroups of $\mathrm{Div}(C)$ and $\mathrm{Pic}(C)$ respectively consisting of the elements of degree zero.
The Jacobian $J(C)$ is then defined to be $\mathrm{Pic}^0(C)$, or in other words, the degree 0 divisors of $C$ modulo linear equivalence.
In fact, it seems to often be defined more abstractly in terms of certain functors, but then the construction is proven to be isomorphic to $\mathrm{Pic}^0(C)$, so we have opted to use that as the definition here.
The group structure is thus the usual group structure in a quotient group: \begin{equation*}
  [D_1] + [D_2] = [D_1 + D_2].
\end{equation*}

\section{Jacobians of genus 2 curves}


\section{Elliptic curves are Jacobians}

\section{The Mordell-Weil group of a Jacobian}

\section{Bibliography}
% https://www.math.ksu.edu/$\sim$ranno/math816/divisors.pdf\\
% Used for learning about divisors and principal divisors.
% (If you want to visit that link, you'll have to change the $\sim$ to a true tilde manually.
% I tried to figure out a way to type it that would get LaTeX to render it correctly but didn't find one.)
\textit{The Arithmetic of Elliptic Curves} by Joseph H. Silverman
\\
Used for understanding the divisor group of a curve.
\\\\
% https://wstein.org/edu/Fall2003/252/lectures/10-15-03/10-15-03-lecture.pdf
% \\
% Used for understanding divisors and principal divisors.
\textit{Algebraic Geometry} by Robin Hartshorne
\\
Used for understanding the formulation of the Jacobian in terms of $\mathrm{Pic}^0(C)$.

\end{document}
